% TRANSLATION UNIT: translate this file only.
% ===========================================================================
%  05 — Metodología
%  The \section title is the canonical (Spanish) chapter name; the body prose
%  is written in ENGLISH and translated to Spanish later.
% ===========================================================================

\section{Metodología}

% DONE: Methodology written in English. The chapter is the honest account of how
% the work was actually done: three accumulated features behind gates, an
% evidence-first discipline with per-chapter claim inventories, the separation
% between machine-verifiable properties and properties that need people, a
% validation design stated before results exist, the threats to validity, and
% what the method deliberately did not do. Prose and tables make only the claims
% inventoried below.

% ---------------------------------------------------------------------------
% CLAIMS INVENTORY — every prose claim is listed here with its evidence.
% The prose below makes only these claims.
% ---------------------------------------------------------------------------
% EVIDENCE: three features accumulate (001 platform → 002 calculator builder and
%   content administration → 003 presentation redesign); each amends the
%   previous, none restarts the product → specs/001-fintcart-platform/plan.md:6-13;
%   specs/002-calculator-builder-content-admin/plan.md:6-11;
%   specs/003-design-system-frontend/plan.md:8-13.
% EVIDENCE: 001's gate is a Constitution Check over the twelve principles, PASS
%   and re-evaluated after Phase 1 → specs/001-fintcart-platform/plan.md:74-93.
% EVIDENCE: 002's gate is the Constitution Check over the design, then over the
%   code, with declared deviations → specs/002-calculator-builder-content-admin/plan.md:107-126,130-175.
% EVIDENCE: 003's gate is the Constitution Check over the design plus a closing
%   re-evaluation (T091) that corrected a real data defect →
%   specs/003-design-system-frontend/plan.md:80-113.
% EVIDENCE: the hybrid working model — Kanban flow with WIP limits plus XP
%   practices inside a document-driven structure — and why not Scrum or pure
%   waterfall → docs/informe/00-plan.md:75-81.
% EVIDENCE: the report's five phases and their gates (skeleton, art direction,
%   calculators, evidence, validation) → docs/informe/00-plan.md:173-181.
% EVIDENCE: the schedule's three cuts (three new seeds, five sessions plus ≥15
%   survey responses, one new diagram) → docs/informe/00-plan.md:333-339.
% EVIDENCE: decision IDs D-26…D-30 exist in both feature 002 and feature 003
%   research files and name different things; every citation is qualified by
%   feature → specs/002-calculator-builder-content-admin/research.md (D-13…D-30);
%   specs/003-design-system-frontend/research.md (D-26…D-32).
% EVIDENCE: the simulator ships ten seed definitions (six general, four
%   Colombian), not the seven the 002 spec still names →
%   services/simulator/src/domain/seeds/mod.rs:144-155.
% EVIDENCE: four roles usuario_final, editor, coordinador_editorial,
%   administrador → services/users/migrations/20260727120000_init_users.up.sql:46;
%   services/users/migrations/20260903090000_roles_administrador.up.sql:15.
% EVIDENCE: every chapter opens with a % EVIDENCE: inventory and its prose may
%   assert only what the inventory lists → docs/informe/chapters/04-marco-referencial.tex:19-33;
%   docs/informe/chapters/06-requerimientos.tex:25-45.
% EVIDENCE: an ad hoc screenshot sweep let two defects ship; an overflow probe
%   and a computed-contrast probe caught them → docs/informe/chapters/11-pruebas.tex:242-266.
% EVIDENCE: a green build does not prove a page renders; a logo mask passed the
%   build and broke the development server → docs/informe/chapters/11-pruebas.tex:223-234,316.
% EVIDENCE: machine-verifiable properties are reported in chapter 11, and
%   human-judgment properties in chapter 13 → docs/informe/chapters/11-pruebas.tex:28;
%   docs/informe/chapters/13-evaluacion.tex:8.
% EVIDENCE: validation is declared before results exist — convenience sample,
%   ≥15 survey responses (target 20), 5–8 moderated sessions, SUS at the end of
%   the session, task success observed → docs/informe/04-instrumentos-validacion.md:282,298,306-309.
% EVIDENCE: the SUS is a perceived-satisfaction instrument, not a learning
%   measure; “preferring is not learning” is the stated limit →
%   docs/informe/04-instrumentos-validacion.md:351,608-616.
% EVIDENCE: honest limits — moderator effect, self-selection, no control group,
%   no significance tests at n ≈ 15 → docs/informe/04-instrumentos-validacion.md:570-601.
% EVIDENCE: the glossary and the prose barrier exist because three authors
%   writing in parallel diverge without a shared, executable gate →
%   docs/informe/00-glosario-y-estilo.md:5; docs/informe/scripts/prose-barrier.mjs:6-10.
% EVIDENCE: no 2PC, no cross-service database reads; saga with compensations
%   only → specs/001-fintcart-platform/plan.md:82,85; .specify/memory/constitution.md:140-142.

% ---------------------------------------------------------------------------
% PROSE
% ---------------------------------------------------------------------------

This chapter describes how the work was actually done, not how a textbook
method would have arranged it. The honest account matters because every result
reported later inherits its limits from the procedure that produced it. The
method has three commitments: the product was built as three increments
accumulated behind gates, every claim rests on evidence that a file, a command,
or a measurement could settle, and machine-verifiable properties were kept
strictly apart from the properties only a person can judge.

\subsection{The shape of the work: three accumulated features, each gated}

This is not three projects. It is one product specified in three features that
accumulated on top of each other, each with its own requirements, its own
acceptance criteria, and a gate that had to pass before the next feature could
begin. Feature \code{001-fintcart-platform} establishes the platform: identity
and access, guided learning, the first calculators, the eight services, and the
constitution of twelve principles that everything later must satisfy. Feature
\code{002-calculator-builder-content-admin} amends it rather than replaces it:
the hard-wired calculators become parameterized seed definitions served by a
formula engine, quizzes become randomized sessions, article bodies become block
documents, and a fourth role is added. Feature
\code{003-design-system-frontend} is different in kind: it is exclusively
presentation, touches nothing outside \code{frontend/}, and is judged by
whether it removes a duplicated style layer rather than by whether it adds one.
Because each feature amended the one before, the numbering of requirements,
acceptance criteria, and decisions is continuous and cumulative, and the
repository is the record of that continuity.

The gate is the same instrument in all three features, and naming it precisely
matters more than naming a process framework. Each feature opens with a
Constitution Check: a table that walks the twelve principles one by one and
states how the design satisfies each of them, and each feature re-runs that
check after Phase 1. Table~\ref{tab:features-gates} records what each gate
demanded.

\begin{table}[htbp]
  \centering
  \begin{tabular}{p{3.2cm}p{4.4cm}p{6.0cm}}
    \hline
    Feature & Gate & What the gate demanded \\
    \hline
    \code{001-fintcart-platform} & Constitution Check, re-run after Phase 1 &
    All twelve principles satisfied by design before any code: gRPC internal
    with REST only at the edge, database-per-service, saga without 2PC,
    decimal-only money, layered services. \\
    \code{002-calculator-builder-content-admin} & Constitution Check over the
    design, then over the code & The same twelve principles, checked against
    the running system rather than the compose file; the seed-regression suite
    must hold before the native calculators retire; declared deviations replace
    silent promises. \\
    \code{003-design-system-frontend} & Constitution Check over the design,
    plus a closing re-evaluation (T091) & No functional behavior may change;
    style debt must reach zero; a monetary figure must never truncate; the
    closing gate corrected a real data defect in the money format. \\
    \hline
  \end{tabular}
  \caption{The three accumulated features and the gate each one had to pass.}
  \label{tab:features-gates}
\end{table}

The first gate demanded that the architecture be constitutional before a line
of code existed, because a violation found after implementation would have
meant rework measured in services, not in files. The second gate mattered most
because it was run over the running system rather than over the compose file:
Redis was checked by listing its actual clients, RabbitMQ by its actual
connections, and the seed-regression suite had to hold before the native
calculator code could retire. The third gate forbade functional change and
demanded that style debt reach zero, and its closing re-evaluation caught a
monetary figure rendered with the wrong thousands separator, a value that
looked right and was another figure entirely.

The same gated pattern governs this document. The writing plan is five phases,
each with a named exit condition: the LaTeX skeleton must compile and match
the institutional template, the art direction must pass its end-to-end and
accessibility gates with zero style debt, the calculators must pass the
seed-regression suite, every figure in the evidence chapters must trace to a
versioned artifact, and validation runs only after the others are closed. A
one-to-two-week deadline then forced three explicit cuts: three new seeds
instead of five, five moderated sessions as the floor, and one new diagram
instead of a redrawn set.

Traceability is not free of collisions, and the method treats them as data
rather than as noise. Decision identifiers overlap between the two research
files: \code{D-26} through \code{D-30} name one thing in feature 002 and a
different thing in feature 003, so throughout this report a decision is always
cited with its feature --- \code{D-27} of feature 003 is the four breakpoints,
and \code{D-27} of feature 002 is the correction of the annuity overflow. Two
authoritative counts are fixed once and reused everywhere: the simulator ships
ten seed definitions (six general, four Colombian), not the seven the 002 spec
still names, and the role model has exactly four roles --- \code{usuario_final},
\code{editor}, \code{coordinador_editorial}, and \code{administrador}.

The working model is the hybrid the plan declares: Kanban flow with
work-in-progress limits, plus the extreme-programming practices of continuous
integration and automated testing, inside a document-driven structure. It is
defended negatively: Scrum would be a fiction for a single developer with no
team, no Product Owner, and no sprints, and pure waterfall is contradicted by
the repository's own history of three features that each amend the previous one
behind a gate.

\subsection{Evidence-first engineering}

The second commitment is a discipline about claims. Every chapter of this
report opens with a \texttt{\% EVIDENCE:} inventory, a list of the specific
claims the prose is allowed to make, each paired with the file and line where a
reader could settle the dispute. The rule is executed rather than recommended:
a chapter may assert only what its inventory lists, and an assertion that
cannot be pointed at a file, a command, or a measurement is reworded or dropped.
This is not documentation theater. It is how three authors writing in parallel
are kept from quietly diverging from the repository, and it is how a reviewer
can tell a measured result from a remembered one.

The discipline changed decisions, not just prose. The clearest case is negative
evidence from feature 003. A first, ad hoc sweep of screenshots was meant to
confirm the redesign by eye, and it let two defects ship: a four-pixel
horizontal overflow on a route the sweep never visited, and a contrast
regression in the token layer, where a value that was true against white was
false against the tinted surface it actually shipped on. Both passed human
inspection, and both were caught only by instruments that assert what nobody
was looking at --- a scroll-width probe and a computed-color contrast check.
The lesson: a review that photographs what it already believes representative
will miss exactly the defect it did not anticipate, and only a measurement that
asserts a property on every target surface can contradict a number that looks
right. Chapter~11 reports both defects in full; this chapter records only what
the episode changed in the method.

\subsection{Two kinds of verification, kept apart}

The method keeps two kinds of verification apart because their failure modes
differ and because conflating them inflates results in one predictable
direction. The first kind is machine-verifiable: a property a command can
assert reproducibly on every run --- unit behavior and numerical edge cases,
end-to-end flows and accessibility selected by role and label, load thresholds,
payload budgets, and style debt. The second kind needs people: whether a task
can be completed without help, how long it takes, and how usable the result
felt. Table~\ref{tab:verification-kinds} assigns each property to its
instrument and its owner chapter.

\begin{table}[htbp]
  \centering
  \begin{tabular}{p{4.4cm}p{5.2cm}p{3.4cm}}
    \hline
    Property & Instrument & Owner chapter \\
    \hline
    Unit behavior and numerical edge cases & unit specs, seed-regression suite & chapter~11 \\
    End-to-end flows and accessibility & Playwright suite selecting by role and label & chapter~11 \\
    Load thresholds & k6 runs on the development host & chapter~11 \\
    Style debt and payload budgets & \code{design-debt.mjs} and bundle checks & chapter~11 \\
    Prose and citation discipline & prose barrier and citation gate & the document itself \\
    Task success, efficacy and efficiency & moderated observation of five tasks & chapter~13 \\
    Perceived usability & SUS administered after the tasks & chapter~13 \\
    Perceived design brief & perception survey & chapter~13 \\
    \hline
  \end{tabular}
  \caption{Properties, their instruments, and the chapter that owns each.}
  \label{tab:verification-kinds}
\end{table}

The separation is not a filing convention; it is the guard against a specific
inflation. A green suite proves that its tests pass and nothing else: it says
nothing about whether a student can complete a task. A favorable SUS score
proves that a participant reported the system as usable and nothing else: it
says nothing about whether any requirement was implemented correctly. If the
two kinds were allowed to borrow each other's authority, chapter~13 could
report a green build as evidence of usability and chapter~11 could report a
favorable survey as evidence of correctness, and neither claim would be earned.
Each chapter therefore declines the other's claims on purpose.

\subsection{Validation design, stated before results exist}

The validation instrument was fixed before any participant was recruited, and
this chapter records it as a commitment rather than as a description of what
happened. The sample is a convenience sample: a minimum of 15 survey responses
with a target of 20, and between 5 and 8 moderated sessions, recruited by
direct invitation from the author's university community. No sampling frame is
claimed, and the instrument states in advance that the sample represents no one
beyond the people who answered. The two instruments measure two different
things. The survey measures declared perception on a five-point anchored scale,
read per item as a distribution, never collapsed into a single average index.
The moderated session measures task success by observation: five tasks, each
with an observable success criterion, a clock, and an escalation ladder that
turns a helped completion into a distinct outcome.

The SUS is used, and it is used for exactly one thing. \parencite{brooke1996sus}
proposed the System Usability Scale in 1996, and that is its role here: a
satisfaction instrument, not a learning measure. Its origin is the only claim
this report makes on the strength of that source, whose full text was
obtainable only as its first page. Its pagination is not cited here: the
chapter DOI resolves to the CRC Press printing, whose record gives pages
207--212, while the usability literature routinely cites 189--194 for the 1996
Taylor \& Francis print, which was not consulted. Rather than pick one of two
unverified numbers, the entry carries no page range and this paragraph carries
the reason. The scale is
administered at the end of the session and reported per session with the
median, never as a percentage of \enquote{usability}.

The honest limits are stated before the data exist because they are properties
of the design, not defects discovered afterward. There is no control group, so
no difference can be attributed to the redesign. The moderator is the author of
the product, so courtesy is a real pressure that the script reduces but cannot
eliminate. The sample self-selects toward people interested in finance or in
the platform. And the rule that governs all of it is negative: preferring the
interface is not evidence of learning. None of these instruments measures
learning; to claim the platform teaches would require a before-and-after
knowledge instrument, a comparison group, and randomization, and none of that
is in this work or simulated here.

\subsection{Threats to validity}

Five threats are named because each one limits what a later chapter may
conclude, and each is a property of the procedure rather than of the results.

First, the load measurement runs on a single development host on a single day;
the figures are a measurement of that host, not a capacity statement about the
deployment. Second, coverage is collected as a diagnostic with no threshold: a
covered line is a line that executed, not a line whose behavior was verified,
so coverage is reported but never used as a quality gate. Third, the end-to-end
suite depends on a running stack; without \code{dev/up} and the seeded data
there is no regression network at all, which is why a broken local environment
reads to the suite as a product defect. Fourth, a passing build does not prove
that a page renders: a logo mask once passed the build and broke the
development server, because the build compiled the template and never painted
it. Visual capture --- actually reading a rendered screen --- is the only gate
that closes that hole, and it is a human judgment. Fifth, the redesign was
measured largely by its own end-to-end suite, which selects by role and
accessible label; that suite proves consistency with the behavior it already
encoded, not that the redesign is good or better. A redesign that passes its
own regression network is self-consistent, not necessarily improved.

\subsection{What the method deliberately did not do}

Three refusals are part of the method and are recorded because a method is
defined as much by what it declines as by what it does.

The system never uses two-phase commit and never reads across service
databases. Multi-service consistency is achieved only through the
orchestrator's sagas with explicit compensations, and the only surface shared
between services is the versioned contract in \code{contracts/}. These refusals
are constitutional, not stylistic: they are the mechanism by which each service
stays independently deployable and each database stays sovereign, and they are
verified in the gates rather than assumed.

The report never lets prose run ahead of a shared vocabulary. Three authors
write chapters in parallel, and without a shared glossary and an executable
barrier they would quietly diverge until the translation stage inherits the
confusion. This is stated plainly because it is a real methodological finding:
the glossary and the prose barrier exist because parallel writing without them
diverges, and the barrier turns the glossary from a recommendation into a gate.
The same discipline applies to citations: a source may be cited by its
existence or by a fact visible in the page actually read, never by its argument
when the full text was not obtained.

Finally, the method never lets a measurement that needs people be replaced by a
measurement that needs a machine, or the reverse. The two are kept apart for
the reason already given: each inflates the other, and a thesis that cannot say
which of its claims are machine-verified and which are human-judged cannot
honestly report either.
